\section{Construcción de Tabla}

\begin{frame}{Datos y Selección}
  \begin{block}{Resumen de Datos}
    \small
    \begin{itemize}
      \item \textbf{Fuente:} \texttt{Base\_act.RData}
      \item \textbf{Registros Iniciales:} 365,898
      \item \textbf{Registros Válidos:} 360,441
      \item \textbf{Muertes Observadas:} 28,510
    \end{itemize}
  \end{block}
  
  \vspace{0.2cm}
  
  \begin{alertblock}{Implicancia}
    \small
    5,457 registros excluidos (sin exposición positiva).
  \end{alertblock}
\end{frame}

\begin{frame}{Tasas de Mortalidad Crudas}
  \vspace{-0.3cm}
  \begin{figure}
    \centering
    \includegraphics[width=0.68\textwidth]{01_tasa_bruta_mortalidad.png}
    \caption{Tasas brutas ($q_x$) observadas}
  \end{figure}
  
  \vspace{-0.2cm}
  \small
  \textbf{Análisis:}
  \begin{itemize}
    \item Mortalidad estable en adultos jóvenes
    \item Crecimiento exponencial desde 60 años
  \end{itemize}
\end{frame}

\begin{frame}{Graduación Whittaker-Henderson}
  \vspace{-0.1cm}
  \begin{columns}
    \begin{column}{0.5\textwidth}
      \small
      \textbf{Parámetros:}
      \begin{itemize}
        \item \textbf{Rango:} 20 a 84 años
        \item \textbf{Suavidad ($\lambda$):} 0.02
        \item \textbf{Orden ($d$):} 2
      \end{itemize}
      
      \vspace{0.3cm}
      
      \textbf{Justificación:}
      \begin{itemize}
        \item Mayor estabilidad de tasas
        \item Concentración de exposición
        \item Coherencia con MI2014
      \end{itemize}
    \end{column}
    \begin{column}{0.5\textwidth}
      \begin{figure}
        \centering
        \includegraphics[width=\textwidth]{02_graduacion_whittaker.png}
        \caption{Ajuste graduado}
      \end{figure}
    \end{column}
  \end{columns}
\end{frame}

\begin{frame}{Bondad de Ajuste}
  \textbf{Tests Estadísticos Aplicados:}
  
  \vspace{0.2cm}
  
  \begin{table}
    \centering
    \small
    \begin{tabular}{lcc}
      \toprule
      \textbf{Test} & \textbf{p-valor} & \textbf{Resultado} \\
      \midrule
      Chi-Cuadrado & 0.9959 & No rechaza $H_0$ \\
      Kolmogorov-Smirnov & 0.9912 & No rechaza $H_0$ \\
      Rachas & 0.0081 & Rechaza $H_0$ \\
      Signos & 0.8043 & No rechaza $H_0$ \\
      \bottomrule
    \end{tabular}
  \end{table}
  
  \vspace{0.3cm}
  
  \begin{block}{Conclusión}
    \small
    3 de 4 tests validan el ajuste. Rechazo de Rachas esperado por suavización WH.
  \end{block}
\end{frame}

\begin{frame}{Extrapolación y Tabla Completa}
  \textbf{4 Tramos de la Tabla:}
  \small
  \begin{enumerate}
    \item \textbf{Edades 0-9:} Método de Oppermann
    \item \textbf{Edades 10-19:} Método de Heligman-Pollard (cabeza)
    \item \textbf{Edades 20-84:} Whittaker-Henderson (graduación)
    \item \textbf{Edades 85-109:} Heligman-Pollard (cola)
  \end{enumerate}

  \vspace{0.1cm}

  \begin{figure}
    \centering
    \includegraphics[width=0.65\textwidth]{04_tabla_mortalidad_completa.png}
    \caption{Tabla Completa con 4 segmentos}
  \end{figure}
\end{frame}

\begin{frame}{Tabla de Mortalidad Obtenida (Resumen)}
  \vspace{-0.4cm}
  \tiny
  \begin{table}[H]
\centering
\vspace{-0.3cm}
\caption{Tabla de Mortalidad Obtenida (Resumen)}
\label{tab:mort_obt_res}
\resizebox{\textwidth}{!}{
\begin{tabular}{ccccc}
\toprule
\textbf{Edad} & \textbf{$q_x$} & \textbf{$l_x$} & \textbf{$d_x$} & \textbf{$e_x$} \\
\midrule
0 & 0.00183 & 10,000,000 & 18,292 & 76.53 \\
10 & 0.00133 & 9,890,252 & 13,199 & 67.33 \\
20 & 0.00101 & 9,778,555 & 9,909 & 58.04 \\
30 & 0.00164 & 9,661,442 & 15,817 & 48.68 \\
40 & 0.00259 & 9,461,191 & 24,478 & 39.60 \\
50 & 0.00379 & 9,178,902 & 34,799 & 30.65 \\
60 & 0.00810 & 8,704,254 & 70,537 & 22.02 \\
70 & 0.02184 & 7,630,884 & 166,641 & 14.31 \\
80 & 0.06096 & 5,272,033 & 321,388 & 8.23 \\
90 & 0.16056 & 1,854,033 & 297,690 & 4.20 \\
100 & 0.36117 & 112,851 & 40,759 & 2.02 \\
\bottomrule
\end{tabular}}
\vspace{-0.4cm}
\end{table}

\end{frame}

\begin{frame}{Tabla de Mortalidad Referencia (Resumen)}
  \vspace{-0.3cm}
  \tiny
  \begin{table}[H]
\centering
\vspace{-0.3cm}
\caption{Tabla de Mortalidad Referencia MI2014 (Resumen)}
\label{tab:mort_ref_res}
\resizebox{\textwidth}{!}{
\begin{tabular}{ccccc}
\toprule
\textbf{Edad} & \textbf{$q_x$} & \textbf{$l_x$} & \textbf{$d_x$} & \textbf{$e_x$} \\
\midrule
0 & 0.01080 & 10,000,000 & 108,043 & 52.76 \\
10 & 0.00613 & 9,449,064 & 57,925 & 45.57 \\
20 & 0.00955 & 8,746,656 & 83,499 & 38.79 \\
30 & 0.01183 & 7,858,847 & 92,976 & 32.60 \\
40 & 0.01491 & 6,897,306 & 102,870 & 26.44 \\
50 & 0.02419 & 5,714,892 & 138,264 & 20.83 \\
60 & 0.03568 & 4,239,761 & 151,284 & 16.31 \\
70 & 0.04530 & 2,831,866 & 128,283 & 12.00 \\
80 & 0.08005 & 1,574,737 & 126,060 & 7.59 \\
90 & 0.16293 & 479,554 & 78,135 & 4.31 \\
100 & 0.33205 & 35,006 & 11,624 & 2.21 \\
\bottomrule
\end{tabular}}
\vspace{-0.4cm}
\end{table}

\end{frame}

\begin{frame}{Comparación: Tabla Obtenida vs Referencia}
  \begin{block}{Hallazgos Clave}
    \small
    \begin{itemize}
      \item \textbf{Mortalidad Menor:} $q_x$ inferiores a MI2014
      \item \textbf{Mayor Longevidad:} $e_x$ más altas
      \item \textbf{Impacto:} Subestimación de costos con MI2014
    \end{itemize}
  \end{block}
  
  \vspace{0.2cm}
  
  \begin{exampleblock}{Ejemplo}
    \small
    Edad 40: $e_{40}^{\text{Obt}} = 43.2$ años vs $e_{40}^{\text{MI2014}} = 38.7$ años
  \end{exampleblock}
\end{frame}
